
% =======================
% Pacotes básicos
% =======================
\usepackage[utf8]{inputenc}
\usepackage[T1]{fontenc}
\usepackage[brazilian]{babel}
\usepackage{amsmath,amssymb}
\usepackage{geometry}
\geometry{
    a4paper, 
    top=3.0cm,    % Aumentei um pouco para caber o cabeçalho de 3 linhas
    bottom=2.5cm, 
    left=2.5cm, 
    right=2.5cm,
    headheight=46pt % Valor corrigido para eliminar o aviso do fancyhdr
}
\usepackage{graphicx}    % Para inserir logo/imagens
\usepackage{adjustbox}   % Para ajustar tabelas
\usepackage{xcolor}      % Cores
\usepackage{hyperref}    % Links clicáveis
\usepackage{minted}      % Código com destaque (compilar com -shell-escape)
\usepackage[most]{tcolorbox}
\usepackage{titlesec}    % Para customizar títulos
\usepackage{fancyhdr}    % Cabeçalho e rodapé
\usepackage{tikz}        % Para desenhar capa com moldura colorida
\usepackage{caption}
\usepackage{subcaption}
\usepackage{ragged2e}
\usepackage{longtable}
\usetikzlibrary{matrix}
\usepackage{comment}
\usetikzlibrary{shapes.symbols, shapes.geometric, arrows.meta, positioning}

% =======================
% Cores personalizadas
% =======================
\definecolor{verdeuni}{RGB}{0,70,40} % Verde do logo UniRV
\definecolor{bgsql}{rgb}{0.95,0.95,0.95}
% FIX: Para o erro do fancyhdr headheight
\setlength{\headheight}{46pt} 
\addtolength{\topmargin}{-20pt}
% =======================
% Cabeçalho e rodapé
% =======================
% --- Configuração de Cabeçalho e Rodapé ---
\pagestyle{fancy}
\fancyhf{}

% Nome do autor no lado esquerdo
\lhead{\color{verdeuni} \footnotesize \textbf{Prof. Me. Sergio S. Novak \\ Prof. Me. Leandro M. P. Sanches \\ Prof. Ronaldo}} 

% Título do guia no lado direito
\rhead{\textbf{\color{verdeuni}Banco de Dados}}

% Número da página no centro do rodapé
\cfoot{\thepage}

% Linha horizontal colorida (opcional, combina com o estilo)
\renewcommand{\headrulewidth}{0.4pt}
\renewcommand{\headrule}{\hbox to\headwidth{\color{verdeuni}\leaders\hrule height \headrulewidth\hfill}}

% =======================
% Sumário automático
% =======================
\setcounter{tocdepth}{2} % mostra até subseções
\hypersetup{
    colorlinks=true,
    linkcolor=verdeuni,
    citecolor=black,
    urlcolor=blue 
}

% =======================
% Estilo de títulos
% =======================
\titleformat{\section}{\Large\bfseries\color{black}}{\thesection}{1em}{} % seções em preto
\titleformat{\subsection}{\large\bfseries\color{verdeuni}}{\thesubsection}{1em}{} % subseções no verde

% =======================
% Caixa para exemplos de SQL
% =======================
\tcbuselibrary{minted}

\newtcblisting{sqlcode}{
    listing engine=minted,
    minted language=sql,
    colback=bgsql,
    colframe=verdeuni,
    breakable,                  % permite quebrar páginas
    listing only,               % só mostra código
    title=SQL,                  % título visível
    fonttitle=\bfseries\color{white},
    coltitle=white,
    colbacktitle=verdeuni
}

\newtcolorbox{diagramER}{
    breakable,
    colback=white,
    colframe=verdeuni,
    title=Diagrama ER,
    fonttitle=\bfseries\color{white},
    coltitle=white,
    colbacktitle=verdeuni,
    boxrule=0.8pt,
    arc=2mm,
    left=4mm,
    right=4mm,
    top=3mm,
    bottom=3mm
}
  
% =======================
% Caixa para observações
% =======================
\newtcolorbox{observacao}{
    colback=green!5,
    colframe=verdeuni,
    boxrule=0.6mm,
    arc=2mm,
    left=4mm,
    right=4mm,
    top=3mm,
    bottom=3mm
}
% =======================
% Definição dos Blocos (FIX para o erro de Environment undefined)
% =======================
\newtcolorbox{block}[1]{%
  colback=green!5,
  colframe=verdeuni,
  coltitle=white,
  title=#1,
  fonttitle=\bfseries,
  boxrule=1pt,
  arc=2mm,
  left=6pt,
  right=6pt,
  top=6pt,
  bottom=6pt
}

\newtcolorbox{quadro_dica}{
    colback=blue!5,
    colframe=blue!50!black,
    title=Dica/Curiosidade,
    fonttitle=\bfseries,
    arc=2mm,
    breakable,
    left=4mm,
    right=4mm,
    top=3mm,
    bottom=3mm
}

\newtcolorbox{danger}{
    colback=red!5,
    colframe=red!75!black,
    title=Cuidado / Erro Comum,
    fonttitle=\bfseries,
    arc=2mm,
    breakable,
    left=4mm,
    right=4mm,
    top=3mm,
    bottom=3mm
}

 


\newtcolorbox{authorcard}{
  colback=white,
  colframe=white,
  boxrule=0pt,
  arc=0pt,          % cantos quadrados
  left=12pt,
  right=12pt,
  top=14pt,
  bottom=14pt,
  height=19.5cm,
  width=140pt,
}
\newlength{\authorphotoheight}
\setlength{\authorphotoheight}{4.5cm}

\newcommand{\authorphoto}[1]{%
\begin{minipage}[c][\authorphotoheight][c]{\linewidth}
\centering
\includegraphics[
width=3.5cm,
height=3.5cm,
keepaspectratio]{#1}
\end{minipage}
}
% =======================
% Dados do documento
% =======================
\title{Banco de Dados}
\author{Me. Sergio Souza Novak} 
\date{\today}

% =======================
% Bibliografia (BibLaTeX)
% =======================
\usepackage[
    style=abnt-numeric,
    backend=biber
]{biblatex}

% Remover cores dos links na bibliografia (Hiperlinks ocultos/pretos)
\AtBeginBibliography{\hypersetup{hidelinks}}

% Link para o arquivo .bib
\addbibresource{bibliografia/normalizacao.bib}
 